% !TEX program = xelatex
\documentclass[11pt,a4paper,fontset=none]{ctexart}
\usepackage[left=22mm,right=22mm,top=23mm,bottom=23mm,headheight=15pt]{geometry}
\usepackage{fontspec}
\setCJKmainfont{Songti.ttc}[Path=/System/Library/Fonts/Supplemental/,FontIndex=6,BoldFont=Songti.ttc,BoldFeatures={FontIndex=1},ItalicFont=Songti.ttc,ItalicFeatures={FontIndex=6}]
\setCJKsansfont{STHeiti Medium.ttc}[Path=/System/Library/Fonts/,FontIndex=1]
\setCJKmonofont{Songti.ttc}[Path=/System/Library/Fonts/Supplemental/,FontIndex=6]
\setmainfont{texgyrepagella}[Extension=.otf,UprightFont=*-regular,BoldFont=*-bold,ItalicFont=*-italic,BoldItalicFont=*-bolditalic]
\setsansfont{texgyreheros}[Extension=.otf,UprightFont=*-regular,BoldFont=*-bold,ItalicFont=*-italic,BoldItalicFont=*-bolditalic]
\setmonofont{lmmono10-regular.otf}[Scale=0.90]
\usepackage{xcolor,graphicx,booktabs,tabularx,longtable,array,multirow}
\usepackage{enumitem,caption,float,fancyhdr,lastpage,needspace}
\usepackage{tikz,pgfplots}
\usetikzlibrary{arrows.meta,positioning,fit,calc,backgrounds}
\pgfplotsset{compat=1.18}
\usepackage[most]{tcolorbox}
\usepackage{hyperref,xurl}
\definecolor{navy}{HTML}{17334D}
\definecolor{teal}{HTML}{008B83}
\definecolor{gold}{HTML}{B67722}
\definecolor{muted}{HTML}{586776}
\definecolor{linegray}{HTML}{DBE3E8}
\definecolor{pale}{HTML}{F1F6F8}
\hypersetup{colorlinks=true,linkcolor=teal,urlcolor=teal,bookmarksnumbered=true,pdftitle={本地大模型技术与应用调研：开放模型、推理服务与业务助手},pdfauthor={MSIM China internship research},pdfsubject={Phase 2 Week 3 technical research; model, serving and business agent applications}}
\urlstyle{same}
\pagestyle{fancy}
\fancyhf{}
\fancyhead[L]{\small\sffamily\color{muted}本地大模型技术与应用研究}
\fancyhead[R]{\small\sffamily\color{muted}PHASE 2 · 2026.09.18}
\fancyfoot[L]{\scriptsize\color{muted}本地大模型技术与应用调研}
\fancyfoot[R]{\small\thepage\ / \pageref*{LastPage}}
\renewcommand{\headrulewidth}{0.3pt}
\renewcommand{\footrulewidth}{0pt}
\setlength{\parskip}{6pt}
\setlength{\parindent}{1.5em}
\linespread{1.22}
\clubpenalty=10000
\widowpenalty=10000
\ctexset{section={format=\Large\bfseries\color{navy},beforeskip=10pt,afterskip=10pt},subsection={format=\large\bfseries\color{navy},beforeskip=9pt,afterskip=5pt},subsubsection={format=\normalsize\bfseries\color{navy}}}
\setcounter{tocdepth}{1}
\setlist[itemize]{leftmargin=1.6em,itemsep=3pt,topsep=3pt}
\setlist[enumerate]{leftmargin=1.8em,itemsep=3pt,topsep=3pt}
\captionsetup{font=small,labelfont={bf,color=navy},skip=6pt}
\renewcommand{\arraystretch}{1.3}
\setlength{\tabcolsep}{5pt}
\newcolumntype{Y}{>{\raggedright\arraybackslash}X}
\newcolumntype{P}[1]{>{\raggedright\arraybackslash}p{#1}}
\newcommand{\src}[1]{\hyperref[ref:#1]{\textcolor{teal}{\mbox{[\csname refnum@#1\endcsname]}}}\allowbreak\hspace{.08em}}
\newcommand{\source}[1]{\par\vspace{1pt}{\footnotesize\color{muted}\raggedright\noindent 依据：#1\par}}
\newcommand{\smallnote}[1]{{\footnotesize\color{muted}#1\par}}
\newtcolorbox{takeaway}[1][]{colback=pale,colframe=teal,boxrule=0pt,borderline west={2pt}{0pt}{teal},sharp corners,left=10pt,right=10pt,top=7pt,bottom=7pt,fonttitle=\bfseries,coltitle=white,colbacktitle=teal,title=#1}
\newtcolorbox{boundary}[1][]{colback=gold!6,colframe=gold!35,boxrule=.4pt,sharp corners,left=8pt,right=8pt,top=5pt,bottom=5pt,fonttitle=\bfseries,coltitle=navy,title=#1}
\newcommand{\referenceentry}[5]{\par\noindent\begin{minipage}{\linewidth}\label{ref:#1}\small\textbf{\color{navy}[\csname refnum@#1\endcsname] #2}\\{\footnotesize #3。#4。\href{#5}{原始出处}}\end{minipage}\par\vspace{3pt}}
\tikzset{flow/.style={draw=linegray,fill=pale,rounded corners=2pt,align=center,minimum height=12mm,text width=25mm,inner sep=4pt,font=\small},flowwide/.style={flow,text width=35mm},arr/.style={-{Stealth[length=2mm]},thick,teal},darr/.style={arr,dashed}}
\newcommand{\flowfive}[5]{\begin{tikzpicture}[node distance=5mm]\node[flow](a){#1};\node[flow,right=of a](b){#2};\node[flow,right=of b](c){#3};\node[flow,right=of c](d){#4};\node[flow,right=of d](e){#5};\draw[arr](a)--(b);\draw[arr](b)--(c);\draw[arr](c)--(d);\draw[arr](d)--(e);\end{tikzpicture}}
\emergencystretch=2em

\newcommand{\product}[1]{\textsf{#1}}

\expandafter\def\csname refnum@TECH-020\endcsname{1}
\expandafter\def\csname refnum@TECH-029\endcsname{2}
\expandafter\def\csname refnum@TECH-037\endcsname{3}
\expandafter\def\csname refnum@TECH-004\endcsname{4}
\expandafter\def\csname refnum@TECH-013\endcsname{5}
\expandafter\def\csname refnum@TECH-014\endcsname{6}
\expandafter\def\csname refnum@TECH-006\endcsname{7}
\expandafter\def\csname refnum@TECH-008\endcsname{8}
\expandafter\def\csname refnum@TECH-016\endcsname{9}
\expandafter\def\csname refnum@TECH-015\endcsname{10}
\expandafter\def\csname refnum@TECH-017\endcsname{11}
\expandafter\def\csname refnum@TECH-012\endcsname{12}
\expandafter\def\csname refnum@TECH-019\endcsname{13}
\expandafter\def\csname refnum@TECH-021\endcsname{14}
\expandafter\def\csname refnum@TECH-022\endcsname{15}
\expandafter\def\csname refnum@TECH-040\endcsname{16}
\expandafter\def\csname refnum@TECH-026\endcsname{17}
\expandafter\def\csname refnum@TECH-009\endcsname{18}
\expandafter\def\csname refnum@TECH-010\endcsname{19}
\expandafter\def\csname refnum@TECH-011\endcsname{20}
\expandafter\def\csname refnum@TECH-054\endcsname{21}
\expandafter\def\csname refnum@TECH-023\endcsname{22}
\expandafter\def\csname refnum@TECH-042\endcsname{23}
\expandafter\def\csname refnum@TECH-024\endcsname{24}
\expandafter\def\csname refnum@TECH-025\endcsname{25}
\expandafter\def\csname refnum@TECH-027\endcsname{26}
\expandafter\def\csname refnum@TECH-055\endcsname{27}
\expandafter\def\csname refnum@TECH-043\endcsname{28}
\expandafter\def\csname refnum@TECH-044\endcsname{29}
\expandafter\def\csname refnum@TECH-045\endcsname{30}
\expandafter\def\csname refnum@TECH-030\endcsname{31}
\expandafter\def\csname refnum@TECH-056\endcsname{32}
\expandafter\def\csname refnum@TECH-049\endcsname{33}
\expandafter\def\csname refnum@TECH-057\endcsname{34}
\expandafter\def\csname refnum@TECH-058\endcsname{35}
\expandafter\def\csname refnum@TECH-053\endcsname{36}
\expandafter\def\csname refnum@TECH-046\endcsname{37}
\expandafter\def\csname refnum@TECH-034\endcsname{38}
\expandafter\def\csname refnum@TECH-033\endcsname{39}
\expandafter\def\csname refnum@TECH-061\endcsname{40}
\expandafter\def\csname refnum@TECH-060\endcsname{41}
\expandafter\def\csname refnum@TECH-059\endcsname{42}
\expandafter\def\csname refnum@TECH-062\endcsname{43}
\expandafter\def\csname refnum@TECH-063\endcsname{44}
\expandafter\def\csname refnum@TECH-035\endcsname{45}
\expandafter\def\csname refnum@TECH-036\endcsname{46}
\expandafter\def\csname refnum@TECH-038\endcsname{47}
\expandafter\def\csname refnum@TECH-039\endcsname{48}
\expandafter\def\csname refnum@TECH-051\endcsname{49}
\expandafter\def\csname refnum@TECH-052\endcsname{50}
\expandafter\def\csname refnum@TECH-032\endcsname{51}
\expandafter\def\csname refnum@TECH-031\endcsname{52}
\expandafter\def\csname refnum@TECH-028\endcsname{53}
\expandafter\def\csname refnum@TECH-041\endcsname{54}
\expandafter\def\csname refnum@TECH-050\endcsname{55}
\expandafter\def\csname refnum@TECH-018\endcsname{56}

\begin{document}
\hypersetup{pageanchor=false}
\begin{titlepage}
\thispagestyle{empty}\setlength{\parindent}{0pt}
\vspace*{16mm}
{\sffamily\large\color{teal}MSIM CHINA · TECHNICAL RESEARCH\par}
\vspace{19mm}
{\Huge\bfseries\color{navy}本地大模型技术\\[5pt]与应用调研\par}
\vspace{9mm}
{\Large\color{muted}开放模型、推理服务与业务助手\par}
\vspace{14mm}\textcolor{teal}{\rule{\linewidth}{2pt}}\vspace{8mm}
\begin{tabularx}{\linewidth}{@{}P{28mm}Y@{}}
研究主题 & 模型能力与许可、Serving framework、Harness与业务应用\\
部署场景 & 专业个人工作站与面向数十人的内部共享服务\\
研究阶段 & Phase 2技术可行性研究，第3周阶段报告\\
资料截止 & 2026年9月18日\\
\end{tabularx}
\vfill
{\large\color{navy}介绍本地模型如何成为可使用的研究工具，\\[4pt]分析模型、服务系统与业务应用之间的关系。\par}
\vspace{16mm}
{\small\color{muted}2026年9月\par}
\end{titlepage}
\hypersetup{pageanchor=true}
\setcounter{page}{1}

\section*{摘要}
\addcontentsline{toc}{section}{摘要}
本地大模型的应用体系由三个相互连接的层次组成：模型承担材料理解、内容生成和工具指令生成；推理服务系统将模型运行在具体硬件上，处理多个请求与上下文；应用及Agent运行环境则把资料、模型和工具组织成用户可以完成的工作。三者共同决定系统的能力、响应速度与使用体验。

当前开放模型已经形成从专业个人工作站到集中计算的多种选择。Qwen和Gemma的中等规模模型提供稠密与混合专家两条路线，兼顾长文阅读、多模态和工具调用；DeepSeek、GLM及Kimi进一步覆盖复杂研究、办公文件和长过程任务。稀疏激活、缓存优化和低精度计算，使这些能力能够对应不同规模的硬件配置。

推理服务为模型提供持续运行的基础。vLLM和SGLang组织多用户请求、上下文缓存与多卡计算，MLX系列发挥Apple统一内存的特点，CPU/GPU异构路线则利用主存和CPU扩展模型容量。新模型通常先获得源码适配，再进入正式发布包；本轮已核对代表模型的部署路径、版本记录和专用构建，为理解安装与维护工作提供具体依据。

非编码应用的主要差别在于工作如何组织。员工工作台支持围绕材料连续提问；固定流程将字段、计算与复核顺序预先确定；通用助手根据目标选择文件和工具操作。以两期财报和一份CSV形成比较表为例，这些形态分别适合探索式分析、重复处理和目标驱动的文件任务，开发框架则支持进一步连接内部系统。

个人工作站便于保留研究环境与本机文件，集中服务可以在多人之间复用模型和资料处理组件。两种方式的经济性可以围绕任务质量、完成时间、资源利用与维护投入展开分析。模型许可中的集团条款、内部使用定义以及应用平台的企业授权，为各类使用方式规定了相应权利与义务。

本报告基于官方模型卡、代码仓库、许可文本和产品文档，介绍主要方案的功能、优势、原理及业务用途，并以参考组合说明模型、推理服务与应用如何共同工作。版本、评测和资料口径集中列于附录。
\clearpage
\tableofcontents
\clearpage
\section{本地模型如何进入基金业务}

基金公司的LLM应用从具体工作开始：研究员希望更快理解材料、查询数据和形成初稿，风控与合规人员需要整理风险信息和审核依据，运营与IT人员则处理文件、异常和代码。模型的语言能力与资料、工具和处理流程结合，形成使用者可以直接操作的业务应用。

例如，研究员要求比较一家公司的数期财报。系统首先取得相关文件，解析并选出与问题有关的内容；模型组织比较方向，必要时调用计算工具核对指标，再形成带来源的分析。模型、资料系统和计算工具分别承担任务，应用负责保存上下文、展示结果并让使用者继续追问。

\begin{table}[H]\small\centering
\caption{基金业务中的工作、模型作用与交付内容}
\begin{tabularx}{\linewidth}{@{}P{22mm}Y Y Y@{}}\toprule
职能 & 具体工作 & LLM承担的工作 & 应用交付的内容\\\midrule
前台投研 & 阅读公告与财报，比较公司和行业变化。 & 理解材料、提取事件、组织证据与解释。 & 比较表、事件记录、研究简报及来源。\\
中台风控与合规 & 整理风险通知，查询制度，预审材料。 & 提取信息、关联条款、标出疑点与修改方向。 & 待复核问题、依据位置与处理记录。\\
后台运营与IT & 处理文件、异常、工单和研究代码。 & 归类信息、形成说明、生成代码或操作步骤。 & 规范化文件、异常说明及经过检查的实现。\\\bottomrule
\end{tabularx}
\end{table}
\smallnote{按公募基金的典型职能归纳。}

从技术组成看，\textbf{模型、Serving framework和Harness分别解决不同问题}。模型接收上下文并产生输出；Serving framework负责模型装载、请求调度、计算和缓存；Harness组织连续的模型调用、工具操作和任务状态。员工工作台或业务应用通常还提供界面、资料管理与用户权限。\src{TECH-020}\src{TECH-029}\src{TECH-037}

\begin{figure}[H]\centering
\begin{tikzpicture}
\node[flow,text width=37mm,minimum height=20mm](a) at (0,0){业务应用与Harness\\界面、任务与工具组织};
\node[flow,text width=37mm,minimum height=20mm](b) at (59mm,0){推理服务系统\\调度、缓存与执行};
\node[flow,text width=33mm,minimum height=20mm](c) at (115mm,0){模型与权重\\理解、生成与推理};
\node[flow,text width=39mm,minimum height=15mm,fill=teal!8](d) at (0,-34mm){资料与研究工具\\检索、数据库、计算};
\draw[arr]([yshift=3mm]a.east)--node[above,font=\scriptsize]{模型请求}([yshift=3mm]b.west);
\draw[arr]([yshift=-3mm]b.west)--node[below,font=\scriptsize]{返回输出}([yshift=-3mm]a.east);
\draw[arr]([yshift=3mm]b.east)--node[above,font=\scriptsize]{执行计算}([yshift=3mm]c.west);
\draw[arr]([yshift=-3mm]c.west)--node[below,font=\scriptsize]{生成结果}([yshift=-3mm]b.east);
\draw[arr]([xshift=-8mm]a.south)--node[left,font=\scriptsize]{工具调用}([xshift=-8mm]d.north);
\draw[arr]([xshift=8mm]d.north)--node[right,font=\scriptsize]{工具结果}([xshift=8mm]a.south);
\end{tikzpicture}
\caption{应用组织模型调用和工具操作；推理服务装载并执行模型。请求与返回分别沿独立箭头传递。}
\end{figure}

这一区分解释了同一模型可以呈现不同的使用体验。接入聊天工作台后，它用于资料阅读与写作；接入固定流程后，它按规定步骤处理公告或待审材料；接入具备文件和计算工具的Agent后，它能够围绕目标持续完成多步工作。应用层的变化也会改变调用次数、上下文长度及服务负载。

本报告围绕专业个人工作站与面向数十人的共享服务展开。前者关注单人任务体验、设备资源和本机工具，后者关注多用户请求、共享资料、用户隔离与运行管理。两种形态可以使用相同模型，也可以由个人应用连接集中推理服务。
\clearpage

\Needspace{8\baselineskip}
\section{开放模型：能力特点与运行资源}

\subsection{模型规模与架构：稠密模型和混合专家}
本次研究覆盖8项开放模型：Qwen和Gemma的4项中等规模模型，面向专业个人工作站及部门服务；DeepSeek、GLM和Kimi的4项大型模型，覆盖更复杂的研究与多步任务。理解这些模型，可以先看两个数字：需要保存多少参数，以及每次计算使用多少参数。

模型名称中的B表示十亿参数，T表示万亿参数。稠密模型在每次计算中使用其主要网络参数；混合专家模型（MoE）保存多组专家，每次选择其中一部分参与计算。例如，Qwen3.6-35B-A3B共包含约35B参数，每次激活约3B。总参数量主要影响权重存储，激活参数量主要影响每次计算，两者共同形成MoE的资源特点。

\subsection{专业个人工作站模型：Qwen与Gemma}
Qwen3.8-27B、Qwen3.6-35B-A3B、Gemma 4 31B和26B A4B都支持长文本及多模态理解。它们为财报阅读、资料问答、表格分析和工具调用提供了专业个人机级别的模型基础，也可以在部门服务器上供多人使用。

\begin{table}[H]\small\centering
\caption{中等规模模型的结构与官方权重体量}\label{tab:personal-models}
\begin{tabularx}{\linewidth}{@{}P{35mm}Y P{25mm}P{20mm}@{}}\toprule
模型 & 模型卡规模口径 & 原生上下文 & 权重体量\\\midrule
Qwen3.8-27B & 语言模型27B，稠密架构。 & 262,144 & 55.56 GB\\
Qwen3.6-35B-A3B & 总35B，每次激活3B。 & 262,144 & 71.90 GB\\
Gemma 4 31B IT & 卡片30.7B，另列视觉编码器。 & 262,144 & 62.55 GB\\
Gemma 4 26B A4B IT & 总25.2B，每次激活3.8B，另列视觉编码器。 & 262,144 & 51.61 GB\\\bottomrule
\end{tabularx}
\source{官方模型卡、配置与固定修订的权重索引。表中均为BF16制品，GB为十进制单位。\src{TECH-004}\src{TECH-013}\src{TECH-014}\src{TECH-006}}
\end{table}

\textbf{Qwen3.8-27B侧重紧凑稠密模型中的综合能力。}它支持文本、图像及视频理解，提供可调节的思考模式，并保留多轮任务的推理上下文。研究员既可以让它快速提取公告字段，也可以增加思考投入，比较材料中的证据、解释差异和组织后续操作。模型卡中的文档理解和工具任务评测，与这类工作直接相关。\src{TECH-004}

\textbf{Qwen3.6-35B-A3B侧重低激活比例的MoE路线。}它以约3B的激活参数参与每次计算，同时保有35B级模型容量，并提供多模态和工具使用能力。这种结构的价值是控制每次调用的计算量，适合持续对话、反复调用工具及多人服务；全部专家权重则约占71.90 GB的BF16存储空间。\src{TECH-013}

\textbf{Gemma 4提供Google的稠密与MoE两条路线。}31B和26B A4B均支持文本、图像、多语言及函数调用，可以连接英文财报、中文资料和图表阅读。31B侧重稠密模型的能力表现，26B A4B通过每次激活约3.8B参数控制计算量。\src{TECH-014}\src{TECH-006}

Google模型卡的128K长文检索测试（MRCR v2，8 needle，平均值）中，31B和26B A4B分别取得66.4\%和44.1\%。该任务要求从长上下文中找回指定信息，31B在这一测试中表现更好。对于同时阅读多份报告的应用，这类成绩提供了模型保持和找回材料信息的比较依据。\src{TECH-014}\src{TECH-006}

\subsection{大型模型：DeepSeek、GLM与Kimi}
大型模型扩大了研究、工具使用和长任务的能力覆盖，同时通过稀疏激活、低精度和缓存设计提高计算效率。本组模型的公开上下文配置均达到约1M，权重制品从数百GB延伸至TB级，对应高内存工作站、多GPU服务器及异构计算等运行形态。

\begin{table}[H]\small\centering
\caption{较大规模模型的结构与当前官方制品}\label{tab:shared-models}
\begin{tabularx}{\linewidth}{@{}P{34mm}Y P{24mm}P{22mm}@{}}\toprule
模型 & 规模及结构 & 制品精度 & 权重体量\\\midrule
DeepSeek-V4.1-Flash & 主干552B，另有196B条件记忆；输入处理/生成阶段激活8B/16B。 & 混合、打包 & 510.29 GB\\
GLM-5.3-Flash & 卡片320B，每次激活18B；混合注意力。 & FP8/BF16 & 328.33 GB\\
GLM-5.3 & 官方制品约753.33B；MoE，沿用GLM-5.2基模。 & FP8/BF16 & 755.62 GB\\
Kimi K3 & 总2.8T，每次激活104B；混合注意力与MoE。 & 4位等混合格式 & 1,560.86 GB\\\bottomrule
\end{tabularx}
\source{官方卡片及权重索引；四项配置均列出约1M上下文上限。权重体量按表列精度统计。\src{TECH-008}\src{TECH-016}\src{TECH-015}\src{TECH-017}}
\end{table}

\textbf{DeepSeek-V4.1-Flash的突出特点是长上下文与Agent执行效率。}它支持图像和文本输入，并将输入处理与逐步生成分为不同计算阶段，分别激活约8B和16B参数。跨层缓存共享、稀疏注意力和低精度缓存减少了保存、读取长上下文的开销，适合反复读取资料和工具结果的研究任务。其指令模型公开评测覆盖自动化、长任务和带工具的图表分析，评测采用最高思考强度。\src{TECH-008}

\textbf{GLM-5.3-Flash将多模态能力与较低计算量结合。}它采用新训练的多模态基模，每次激活约18B参数，通过线性注意力与稀疏注意力处理长材料。文本、图像和工具调用可以组成连续任务，例如读取财报页面、提取表格，再调用计算工具形成比较结果。当前权重约328.33 GB，是本组中体量较小的一项；官方同时给出了GPU和CPU/GPU异构部署路径。\src{TECH-016}

\textbf{GLM-5.3的重点是复杂工具任务与长过程执行。}它沿用GLM-5.2基模，通过后训练增强代码、自动化和多步任务能力。发布方在Toolathlon Verified中报告73.0，采用官方评测服务的三次独立运行平均单次成功率口径。对于研究系统，这类能力可以用于组织数据查询、计算脚本、文件处理及连续工具操作。\src{TECH-015}

\textbf{Kimi K3的特点是研究、办公和多模态任务覆盖较完整。}模型卡包括办公PDF、电子表格、公司金融与金融Agent评测，展示了从读取材料到操作工具、交付文件的能力。卡片引用Vals AI的CorpFin v2与Finance Agent v2成绩分别为71.6和54.4，对应最高思考强度模型列。它采用2.8T规模的MoE，每次激活104B，并以混合低精度权重控制存储；当前约1.56 TB的制品对应更大规模的集中资源配置。\src{TECH-017}

\Needspace{16\baselineskip}
\subsection{模型运行中的显存与内存需求}\label{sec:memory}
运行一个模型，内存主要用于三件事：\textbf{保存模型权重、保留当前任务的上下文、支持计算过程}。普通GPU服务器主要使用显存；Apple统一内存和CPU/GPU异构路线则让系统内存参与模型运行。图\ref{fig:memory}给出这三部分的关系。

\begin{figure}[H]\centering
\begin{tikzpicture}[node distance=5mm]
\node[flow,text width=32mm,minimum height=18mm](a){模型权重\\保存模型学到的参数};
\node[flow,right=of a,text width=32mm,minimum height=18mm](b){任务上下文\\保存已读材料的状态};
\node[flow,right=of b,text width=32mm,minimum height=18mm](c){计算空间\\保存计算中的临时数据};
\node[below=9mm of b,font=\small\bfseries\color{navy}](d){三部分共同占用显存或内存};
\draw[arr](a.south)--(d.north west);\draw[arr](b)--(d);\draw[arr](c.south)--(d.north east);
\end{tikzpicture}
\caption{模型运行的三类主要内存用途。}\label{fig:memory}
\end{figure}

\textbf{权重决定基本装载规模。}Qwen3.8-27B的BF16权重约55.56 GB，Qwen3.6-35B-A3B约71.90 GB。后者每次只选择部分专家计算，但内存中保存的是全部专家。因此，选择模型时，总体量回答“需要多大的存储空间”，激活量回答“每次调用要做多少计算”。对于48 GB显存配置，可通过降低精度、跨设备分配或CPU/GPU协作安排这两款模型。

\textbf{上下文占用随任务长度和同时运行的请求变化。}一个短问题只包含少量输入；两期财报比较则包含长篇材料、提问和工具结果。推理系统把已经处理过的信息保存为上下文状态，其中常见形式是键值缓存（KV cache），以便继续生成和追问。多人共用一份模型权重时，各个请求还会产生相应的上下文状态，共同占用剩余空间。

\textbf{计算过程也需要工作空间。}矩阵计算、图像处理和加速执行会使用临时数据与缓冲区。由此，设备配置先容纳权重，再为预期的材料长度、同时处理的请求及计算过程分配余量。

\textbf{量化通过减少数值位宽节省空间。}例如，BF16使用16位浮点表示，FP8和4位格式进一步压缩相应权重或缓存。量化让较大模型进入更紧凑的设备，也为更长上下文或更多请求留出空间；其计算效率由对应硬件和推理内核共同发挥。

\Needspace{8\baselineskip}
\section{许可：模型权重与机构使用方式}

\subsection{开放权重对应不同的授权条件}
模型能否取得权重、采用哪种许可，以及机构怎样使用，是相互关联的三个问题。公开权重提供了本地运行和进一步研究的条件；具体的修改、分发、商业使用和附加义务，则由相应许可决定。

本轮Qwen与Gemma候选采用Apache 2.0，DeepSeek-V4.1-Flash和GLM-5.3-Flash采用MIT。两类许可允许相应的使用、修改和分发，并规定声明保留等条件。Apache 2.0还明确列出特定范围的专利许可、分发时的修改标记和NOTICE处理，其法律主体定义涉及控制、被控制或共同控制的实体。\src{TECH-004}\src{TECH-013}\src{TECH-012}\src{TECH-019}\src{TECH-008}\src{TECH-016}

GLM-5.3与Kimi K3采用各自的自定义许可。它们的重要区别集中在关联方、模型即服务业务和内部使用范围。表\ref{tab:license}概括本轮核对的条件。

\begin{table}[H]\small\centering
\caption{主要模型许可及使用条件}\label{tab:license}
\begin{tabularx}{\linewidth}{@{}P{30mm}P{29mm}Y@{}}\toprule
模型组 & 许可 & 与机构使用相关的内容\\\midrule
Qwen两项、Gemma两项 & Apache 2.0 & 版权及特定专利授权；分发时保留许可与声明，标记修改，按条件处理NOTICE。\\
DeepSeek-V4.1-Flash、GLM-5.3-Flash & MIT & 允许相应商业使用、修改和分发；保留版权及许可声明。\\
GLM-5.3 & GLM-5.3 License & 关联方经营MaaS与集团连续12个月收入门槛共同触发安全审查要求。\\
Kimi K3 & Kimi K3 License & MaaS协议、商业产品署名及内部使用豁免分别规定。\\\bottomrule
\end{tabularx}
\source{各模型固定修订的LICENSE、Google官方Gemma 4许可页及Apache全文。\src{TECH-015}\src{TECH-017}}
\end{table}

\subsection{集团主体与内部使用的差别}
GLM-5.3的附加条件同时考察被许可方或关联方是否经营模型即服务（MaaS），以及集团合计收入是否超过门槛；满足联合条件时，商业使用前涉及Z.AI安全审查。Flash采用MIT，适用标准MIT授权条件。主模型的判断单位则同时涉及关联实体、集团收入和实际经营的服务内容。\src{TECH-015}\src{TECH-016}

Kimi K3也有MaaS及商业产品条件，但另行规定内部使用豁免；其定义要求软件、输出及底层能力均不向第三方提供。GLM主模型没有同样的内部使用豁免。纯内部研究、集团跨实体共享与向客户交付内容，分别对应主体认定、第三方范围及输出提供方式等适用问题。具体金额、门槛和条款位置集中列于附录\ref{app:model-license}。\src{TECH-017}

\subsection{模型来源、软件组件与内部采用}
Google的Gemma与Qwen、DeepSeek、GLM及Kimi构成不同发布方的模型路线。机构使用时，可以分别识别权重发布者、服务供给主体和数据处理位置：本地部署由内部设备承担推理，使用托管服务则涉及相应服务商及其运行环境。

应用系统还包含推理引擎、资料解析、嵌入、工具、界面和第三方扩展。这些组件各自适用相应许可。例如，一款MIT模型可以接入带商业条件的应用平台，形成模型和平台分别授权的组合。产品中的企业身份管理、支持服务和品牌修改也可能单独授权。软件许可与产品边界汇总见附录\ref{app:software-license}。


\Needspace{8\baselineskip}
\section{推理服务：模型怎样持续回应使用者}

\subsection{从单次生成到多用户服务}
Serving framework将模型权重装载到设备上，并持续处理应用发来的请求。对使用者而言，它主要影响开始回答前的等待、持续生成的速度，以及多人同时使用时的响应表现。对运维人员而言，它组织计算、缓存、调度和资源使用。

推理过程通常包含输入处理（prefill）与逐步生成（decode）。阅读长财报时，输入处理会占用较多计算；生成较长分析时，后续decode持续使用模型。一个研究Agent在取得检索结果或计算返回值后，通常还会继续调用模型，因此其服务负载包含多轮输入处理和生成。

\begin{figure}[H]\centering
\begin{tikzpicture}[node distance=6mm]
\node[flow,text width=30mm](a){应用提交上下文\\问题、资料、工具结果};
\node[flow,right=of a,text width=29mm](b){输入处理\\读取现有上下文};
\node[flow,right=of b,text width=29mm](c){逐步生成\\答案或工具指令};
\node[flow,below=10mm of b,text width=37mm,fill=teal!8](d){应用继续处理\\展示、调用工具或结束任务};
\draw[arr](a)--(b);\draw[arr](b)--(c);\draw[arr](c.south)|-(d.east);
\draw[darr](d.west)-|node[pos=.22,below,font=\scriptsize]{新的工具结果}(a.south);
\end{tikzpicture}
\caption{业务任务中的模型调用。多步任务会多次经过输入处理和生成，完整体验还包含工具执行时间。}
\end{figure}

连续批处理可以让多个请求共用计算批次，前缀缓存可以复用重复上下文，多卡并行则将模型或请求分配到多个设备。这些机制解释了推理服务的主要价值：在模型能力基本不变的情况下，改善硬件利用与用户响应。\src{TECH-020}\src{TECH-021}

\subsection{共享GPU服务：vLLM、SGLang与TokenSpeed}
vLLM与SGLang均围绕持续模型服务构建，提供批处理、缓存、模型适配、结构化输出和多卡执行能力。它们通过集中调度多个请求提高设备利用率，并用缓存与并行执行改善长材料、多轮任务的响应。

\begin{table}[H]\small\centering
\caption{共享GPU服务中的主要框架}
\begin{tabularx}{\linewidth}{@{}P{29mm}Y Y@{}}\toprule
项目 & 做什么 & 主要技术特点与价值\\\midrule
vLLM & 将模型提供为持续运行的HTTP服务。 & 分页式注意力内存管理、连续批处理、缓存和多后端生态，适合共享模型服务。\\
SGLang & 组织高吞吐推理、上下文缓存和分布式执行。 & RadixAttention前缀复用、调度及多种并行能力，适合长上下文与大型MoE服务比较。\\
TokenSpeed & 围绕Agent负载优化推理执行。 & 统一管理请求和缓存，替换针对设备优化的计算组件，关注多轮Agent调用效率。\\
TensorRT-LLM & 面向NVIDIA设备优化模型推理。 & 结合目标设备、精度与模型实现的专用路线，适合固定组合的深入比较。\\\bottomrule
\end{tabularx}
\source{各项目官方README及服务文档。\src{TECH-020}\src{TECH-021}\src{TECH-022}\src{TECH-040}}
\end{table}

vLLM使用PagedAttention等机制管理注意力缓存，并提供多种模型格式、推理解析器和API。其服务接口覆盖Chat Completions与Responses等协议，使应用能够通过统一接口调用不同模型。\src{TECH-020}

SGLang的RadixAttention利用多个请求之间可复用的前缀减少重复计算，并支持分块prefill、连续批处理和张量、流水线、专家、数据等并行方式。对于重复读取同一资料或在相近上下文中继续工作的问题，缓存命中和调度方式会影响实际收益。当前SGLang也可作为KTransformers异构执行路线的服务层。\src{TECH-021}\src{TECH-026}

TokenSpeed将请求调度、缓存管理与具体模型计算分开，便于围绕Agent反复调用的特点优化执行。其后续源码已提供Qwen3.8-27B、GLM-5.3及Flash、Kimi K3的部署配置，形成面向当前模型和特定设备的性能优化路线。\src{TECH-022}

\begin{table}[H]\small\centering
\caption{代表模型的部署支持证据与构建条件}\label{tab:build-support}
\begin{tabularx}{\linewidth}{@{}P{28mm}Y Y@{}}\toprule
模型与引擎 & 本轮取得的具体证据 & 对采用条件的含义\\\midrule
Qwen3.6 / vLLM、SGLang & 模型卡列出vLLM至少0.19.0、SGLang至少0.5.10；vLLM 0.29.0已有架构入口。 & 可沿已有正式版本配置模型，配套选择精度与工具解析器。\\
Qwen3.8 / vLLM & 官方部署配置记录具体设备与测试构建，部分为开发版本或专用镜像。 & 可按上游记录的设备、精度和镜像组合复现部署。\\
DeepSeek V4.1 / vLLM & 初始支持于9月11日合并，晚于9月9日的v0.29.0发布。 & 本轮固定正式版与后续源码不同，复现需包含相应适配的构建。\\
GLM-5.3-Flash / vLLM & 后续源码已有新架构入口，本轮v0.29.0注册表未见该入口。 & 部署路线跟随后续源码适配，并配套模型格式与工具解析。\\
新模型 / TokenSpeed & v0.1.0发布于7月24日；上述模型配置来自后续源码。 & 采用包含相应模型配置的后续源码，并固定设备与修订。\\\bottomrule
\end{tabularx}
\source{模型卡、正式版注册表、后续源码、合并及发布记录。\src{TECH-013}\src{TECH-009}\src{TECH-010}\src{TECH-011}\src{TECH-054}\src{TECH-020}\src{TECH-022}}
\end{table}

成熟引擎提供调度、缓存和并行等共用基础，新模型通过专用适配接入这些能力。正式发布包便于统一安装与升级，后续源码能够更早使用新模型；二者主要对应不同的版本固定和维护方式。

\subsection{Apple Silicon：文本、多模态与Metal路线}
MLX是Apple Silicon上的数组计算框架，MLX-LM在其上提供文本模型装载、量化、生成和微调，并通过服务接口支持Agent应用。Apple在WWDC26演示了MLX-LM Server连接OpenCode、连续批处理和多Mac分布式推理。\src{TECH-023}\src{TECH-042}

MLX-VLM是基于MLX的社区多模态项目，扩展到图像、音视频及相应模型服务。其当前文档包含Qwen3.8和Gemma 4的示例，并提供Chat Completions、Responses及结构化输出等能力。对于财报图表和扫描材料，这一方向将视觉理解与文本任务连接起来。\src{TECH-024}

llama.cpp则提供Metal后端与GGUF格式，可在Mac及其他平台上运行。它的价值在于跨平台格式与执行路径较灵活，既可用于个人服务，也支持CPU/GPU混合。Ollama、LM Studio等应用入口还可以在底层引擎之上承担模型管理与使用界面。\src{TECH-025}\src{TECH-042}

Apple统一内存使CPU与GPU访问同一内存体系，模型、缓存、系统和其他应用共享总容量。较大的统一内存可以容纳更大模型；多Mac运行进一步通过分片和互联组织多台设备，扩展总资源。

\subsection{CPU/GPU异构：利用主存与CPU参与大型模型计算}
CPU/GPU混合执行的共同思路，是将模型的部分权重或计算放到CPU一侧，以主存容量和CPU计算配合GPU。可以按模型层或专家分配工作，实际效果受主存带宽、CPU计算能力、处理器与内存的连接方式，以及CPU和GPU之间的传输影响。

KTransformers当前以kt-kernel提供CPU优化内核和异构MoE执行，可与SGLang服务层集成。其GLM-5.3-Flash教程列出CPU专家计算、GPU专家配置和分层prefill，并给出特定GPU与CPU内核条件。该教程中的FP8权重约306 GiB，系统内存预留条件达到数百GB，说明这一路线也涉及专业级内存与主机配置。\src{TECH-026}

llama.cpp同时支持CPU、全GPU和混合执行，提供较通用的模型装载与服务方式。其独立维护分支ik\_llama.cpp增加特定量化、内核和并行优化。llama.cpp提供通用执行基础，ik分支则围绕特定量化和内核进一步优化性能；后者的维护内容还包括跟进上游模型与工具格式。\src{TECH-025}\src{TECH-027}

\begin{table}[H]\small\centering
\caption{个人及异构推理路线的差异}
\begin{tabularx}{\linewidth}{@{}P{29mm}Y Y@{}}\toprule
路线 & 优势所在 & 影响使用体验的条件\\\midrule
MLX-LM / MLX-VLM & 利用Apple Silicon与统一内存，分别覆盖文本和多模态。 & 模型转换、精度、上下文、API和实际Mac配置。\\
llama.cpp & GGUF、多后端和灵活装载，兼顾Metal及CPU/GPU混合。 & 量化质量、内存带宽、权重放置及并发。\\
KTransformers & CPU专家计算与GPU协作，扩大大型MoE的实现方式。 & CPU计算、主存连接、专家放置与服务组合。\\
ik\_llama.cpp & 针对特定量化与设备的优化实现。 & 对应模型收益、构建和与上游的维护差异。\\\bottomrule
\end{tabularx}
\end{table}

\subsection{速度分析与业务负载}
员工交互、批量公告处理和多步研究Agent，对服务提出不同要求。交互问答关注等待时间，批处理关注一定时间内完成的工作量，多步Agent还包含工具调用与反复prefill。用户人数、活跃会话数和同时生成的请求数由此形成不同的容量口径。

vLLM的benchmark工具能够记录请求速率、时延分位数和goodput。这里的goodput表示满足指定时延目标的请求吞吐；金融业务还需要加入答案与产物质量，才能进一步衡量单位时间内完成的合格任务。\src{TECH-054}

SGLang的调优文档将离线批处理作为部分参数说明的明确场景，并讨论缓存池、调度保守程度和prefill大小。离线任务可以通过较长队列和较大批次提高吞吐；员工交互更关注及时响应，对应较短等待和峰值请求管理。\src{TECH-055}

从经济性看，任务质量、响应时间和吞吐共同影响业务产出。系统在相同时间内交付更多可用结果，并减少重新生成和人工修订，就能提高完整任务的处理效率；这些因素连接了推理服务与下一章的应用层。

\Needspace{8\baselineskip}
\section{Harness与业务应用：同事最终使用什么}

\subsection{从模型接口到员工工作台}
非编码工作同样需要模型之外的能力。读取财报涉及文件解析与检索，分析表格涉及计算，整理研究材料涉及引用、文件生成和任务记录。Harness将模型与这些工具组织起来，完整应用再提供同事能够直接使用的界面。

因此，应用层可以按最终工作形态理解：员工工作台提供日常入口，资料研究平台连接知识与证据，流程平台处理重复工作，通用助手围绕目标持续行动，开发框架则用于构建更贴合内部系统的应用。

\begin{table}[H]\small\centering
\caption{AI应用的主要形态与业务价值}\label{tab:application-types}
\begin{tabularx}{\linewidth}{@{}P{28mm}Y Y@{}}\toprule
应用形态 & 代表项目 & 使用者得到的能力\\\midrule
员工AI工作台 & Open WebUI、LibreChat、AnythingLLM & 在统一入口聊天、上传文件、使用资料与专用助手。\\
资料研究平台 & Onyx、RAGFlow & 在大量资料中查找证据、比较信息、形成可追溯回答。\\
固定流程平台 & Dify、n8n、Flowise & 把一次有效操作固化为可以重复运行的处理流程。\\
通用任务助手 & Goose、OpenClaw & 接收一个目标，组合文件和工具操作，持续交付结果。\\
应用开发框架 & LangGraph、Microsoft Agent Framework、Agno、CrewAI & 由工程团队定义业务步骤、工具、状态和人工处理。\\\bottomrule
\end{tabularx}
\source{根据官方产品定位与功能归纳，类别可重叠。\src{TECH-043}\src{TECH-044}\src{TECH-045}\src{TECH-030}\src{TECH-047}\src{TECH-049}}
\end{table}

\subsection{贯穿示例：两期财报与一份CSV形成比较表}\label{sec:case}
以下使用虚构公司和数据说明工作过程。输入包括云川公司2024年、2025年财报各一份，以及按年、季度记录营业收入的CSV。两份财报均以人民币亿元列示持续经营业务，营业收入分别为120、138，营业成本分别为84、96.6；CSV中两年的季度收入分别为28、30、29、33和32、34、35、37。

任务是形成带期间、单位和来源的比较表，并给出三条说明。模型识别财报字段与注释，计算工具完成单位对齐、加总和比率计算，再由应用组织引用与结果文件。表\ref{tab:case-output}展示所需产物。

\begin{table}[H]\small\centering
\caption{财报比较任务的说明性输出}\label{tab:case-output}
\begin{tabularx}{\linewidth}{@{}P{30mm}P{20mm}P{20mm}Y@{}}\toprule
指标 & 2024年 & 2025年 & 来源及计算依据\\\midrule
营业收入（亿元） & 120 & 138 & 各期财报第12、14页；CSV季度加总分别一致。\\
营业成本（亿元） & 84 & 96.6 & 各期财报同页，币种、单位及业务范围一致。\\
毛利率 & 30.0\% & 30.0\% & （营业收入减营业成本）除以营业收入。\\\bottomrule
\end{tabularx}
\smallnote{示例设定：公司、数字和页码均为虚构；金额统一为人民币亿元。}
\end{table}

配套说明为：营业收入增长15\%；营业成本同比例增长，毛利率保持30\%；两年CSV的季度加总分别为120和138亿元，与财报一致。这个任务把材料理解、确定性计算和来源核对组成一个完整产物。下文几种产品形态的差别，主要体现在谁来安排这些步骤。

\subsection{员工工作台：聊天、文件与可复用助手}
\subsubsection{Open WebUI}
Open WebUI提供自托管的多模型工作台。员工通过浏览器选择模型、上传资料、使用知识库或调用已经配置的工具，后台连接则由统一平台管理。其官方功能包括用户组、角色权限、模型配置和可扩展工具。\src{TECH-043}

这一形态的优势在于把分散的模型使用方式集中到一个入口。研究员可以围绕公开财报连续提问，运营人员可以查询允许访问的制度资料，团队还可以保存面向不同工作的专用配置。模型服务与使用界面分离，也便于后台更换模型。

\subsubsection{LibreChat与AnythingLLM}
LibreChat将多模型聊天、专用Agent、文件搜索和工具操作结合。用户可以为一个助手保存指令、资料和工具，并将其共享给同事。其当前文档还覆盖文件上传、代码执行与生成结果下载。\src{TECH-044}

在上述比较任务中，使用者可以先上传三份文件，要求列出指标，再追问增长率、要求核对CSV，最后下载结果；任务方向由连续提问逐步确定。LibreChat可将文件搜索、计算与下载接在同一会话中，其优势是保留探索和修正空间。代码执行通过可自托管的计算服务提供；持久会话和附接工作环境则是当前标为实验性的新功能。\src{TECH-044}（Agents、Code Interpreter文档）

AnythingLLM将模型连接、资料导入、文档聊天与Agent能力整合在一起，提供桌面和自托管方式。围绕行业或项目建立资料工作空间后，使用者可以持续整理材料、查询信息和形成草稿。其优势是个人资料与模型入口比较集中；当前README将多人及权限能力对应到Docker路线，桌面、团队部署和Pro产品提供的管理功能不同。\src{TECH-053}

\subsection{资料研究平台：让回答能够回到证据}
\subsubsection{Onyx：连接企业资料的搜索与研究助手}
Onyx连接已有信息源，提供搜索、问答、专用Agent、多步研究和文件产物。它既包含面向员工的使用入口，也包含完整部署所需的资料索引与检索组件。\src{TECH-045}

对于基金研究，这种产品形态可以围绕公司或行业持续查找依据。例如，系统取得允许访问的材料，比较不同时间的观点，再组织为带证据的研究简报。其主要价值在于将资料发现与研究问题放在同一工作过程，减少逐个平台寻找信息的时间。

企业资料的访问范围会随人员和文档变化，因此连接器和权限继承是这类系统的重要组成。Onyx提供社区核心和企业功能，企业身份与权限等能力按相应产品版本提供，企业代码目录采用单独许可。\src{TECH-045}（LICENSE及产品版本页面）

\subsubsection{RAGFlow：解析复杂文档并组织可追溯回答}
RAGFlow重点处理资料解析、切分、索引和基于证据的回答，并将检索与Agent流程结合。对于长PDF、跨页表格或包含注释的财报，资料如何被转换为模型上下文，会直接影响最终答案。\src{TECH-046}

RAG流程可以概括为：系统先将文档整理成可检索资料；使用者提问后，检索组件选取有关片段；模型据此形成回答，应用保留来源位置供核查。RAGFlow的特色在于资料处理与引用检查，而Open WebUI、Dify等产品也可以利用相同机制提供知识问答。

\begin{figure}[H]\centering
\flowfive{原始资料\\财报、公告、制度}{文档处理\\解析、切分、索引}{查找证据\\问题与相关片段}{模型回答\\归纳、比较、解释}{结果核查\\引用与原文位置}
\caption{资料问答的共同流程。不同产品的差异主要体现在文档处理、资料连接、检索与用户核查方式。}
\end{figure}

Onyx与RAGFlow的侧重点由此形成区别：前者强调连接企业知识并提供研究入口，后者突出文档成为高质量上下文的过程。在前述例子中，如果财报尚未收集，Onyx类平台侧重从已有资料源找到正确期间的报告；当PDF已经取得，RAGFlow类平台侧重保留跨页表格、单位与脚注，便于回答回到原文。本例的三份指定材料可直接解析后交给模型；资料持续积累时，索引和检索使已有材料能够重复利用。

\subsection{固定流程平台：把重复工作变成应用}
\subsubsection{Dify：通过可视化流程提供部门应用}
Dify将模型、资料、工具和判断条件放在可视化工作流中，可以形成聊天助手或固定业务应用。其价值是让处理过程可见，并将模型配置、知识库、应用入口及运行记录集中起来，便于业务人员与工程团队讨论同一套流程。\src{TECH-030}

同一财报比较任务可以预先固定为：接收两份报告和CSV，提取年度、币种、单位、收入与成本，程序计算增长率和毛利率，并核对季度加总。若期间或单位无法对齐，流程将疑点和出处交给人员确认，再保存比较表。Dify提供搭建与运行这些步骤的平台，应用配置明确字段、核验规则、导出和复核入口，使使用者可以重复运行同一方法。

\begin{figure}[H]\centering
\flowfive{接收材料\\两期财报与CSV}{模型提取\\期间、单位与字段}{程序计算\\加总、增长与毛利率}{人员确认\\疑点与出处}{保存产物\\比较表与三条说明}
\caption{财报比较的固定流程设计：从接收材料、提取和计算，到人员确认及结果保存。}
\end{figure}

Dify自托管还包含数据库、缓存、插件、后台任务进程及隔离的计算环境。可视化搭建整合界面与流程，后台组件负责状态保存、任务执行和计算隔离；多工作空间和品牌条件见附录\ref{app:software-license}。\src{TECH-030}

\subsubsection{n8n与Flowise：业务连接和可视化搭建}
n8n侧重连接现有业务系统，通过触发器把输入、规则、模型处理和后续动作串起来。例如，新资料进入后可以依次完成归类、整理和归档，模型参与其中需要语言理解的环节。它的优势主要体现在跨系统自动化及已有连接器。\src{TECH-034}

Flowise通过Chatflow与Agentflow展示模型、资料和工具之间的关系，便于试验不同处理流程。与Dify相似，它为搭建应用提供可视化入口，团队管理和企业功能按版本与许可提供，汇总见附录\ref{app:software-license}。\src{TECH-033}\src{TECH-034}

这些平台改变的是工作组织方式。它们把重复的人工操作转成可运行、可检查的过程，模型负责语言理解，资料、规则和工具提供处理依据与执行能力。

\subsection{通用任务助手：围绕一个目标持续完成工作}
\subsubsection{Goose与OpenClaw}
Goose提供桌面应用、CLI和API，通过模型与工具完成研究、写作、自动化和数据分析。其当前官方项目为aaif-goose/goose，具有自定义模型服务和扩展能力，并以Recipes保存可复用的指令、工具及设置。\src{TECH-047}

对于同一比较任务，研究员可以一次提出“读取这两期财报和CSV，生成比较表及三条说明”。Goose依据目标选择读取、计算和写文件的步骤；这些能力通过配置的文件、解析和执行工具提供。处理顺序由助手在任务中决定，便于临时增加分析问题或更换产物格式。桌面入口提供交互界面，Recipes保存有效的方法，扩展机制连接所需的文件和业务工具。

OpenClaw更强调持续个人任务。它将会话、记忆、工具和不同入口连接到自托管Gateway，使关注主题和工作背景可以延续到多次交互。资料整理、会议准备和持续跟踪因此可以围绕同一个助手开展。\src{TECH-048}

两者的使用形态各有重点：Goose提供面向电脑任务的通用操作入口，OpenClaw强调持续状态与多渠道访问。OpenClaw官方以一个Gateway对应一个信任边界；当不同人员需要严格隔离时，其个人助手配置与企业共享平台存在不同的管理要求。\src{TECH-048}（Trust model）

\subsection{开发框架：把业务逻辑落实为专用系统}
现成应用提供使用入口，开发框架则让团队进一步掌握处理步骤、工具和任务状态。对于具有特定数据接口、审批过程或输出格式的工作，框架决定了这些规则怎样进入实际应用。

\begin{table}[H]\small\centering
\caption{开发框架的主要作用与优点}
\begin{tabularx}{\linewidth}{@{}P{35mm}Y Y@{}}\toprule
项目 & 做什么 & 主要价值\\\midrule
LangGraph & 组织步骤、分支、任务状态与人工中断。 & 过程清楚，可把固定流程与开放研究结合。\\
Microsoft Agent Framework & 提供Agent与多步流程的统一开发基础。 & 连接多模型、工具与企业应用的路线较完整。\\
Agno & 组合Agent、协作团队、流程及运行管理。 & 将知识、记忆和助手服务组织在同一产品路线。\\
CrewAI & 按职责组织Agent，并用Flow控制过程。 & 资料搜集、核对和写作等分工较直观。\\
Haystack & 将检索、排序、过滤和生成组成处理流程。 & 资料处理和答案依据较容易分环节检查。\\
LlamaIndex / LangChain & 提供文档、索引、模型及工具连接组件。 & 复用现有集成，减少重复建设。\\\bottomrule
\end{tabularx}
\source{官方README与概念文档。\src{TECH-029}\src{TECH-049}\src{TECH-051}\src{TECH-052}\src{TECH-032}\src{TECH-031}\src{TECH-028}}
\end{table}

LangGraph组织状态、分支和人工中断，LangChain提供模型与工具连接。Deep Agents进一步组合规划、文件工具、上下文管理和子任务，让工程团队复用多步助手的常见能力。这些能力为开发专用研究助手提供了现成组件。\src{TECH-028}\src{TECH-029}\src{TECH-041}

DeepSeek Harness同样提供可组合的模型、工具、会话和执行循环，适合定制任务组织方式。当前项目处于开发者预览阶段，重点提供可组合的执行机制，工程团队在其上建设业务入口、数据连接和运行管理。\src{TECH-037}

Microsoft Agent Framework提供.NET和Python等开发路线，包含工具、工作流、人工处理和多模型服务连接。微软官方说明其1.0已于2026年4月正式发布；AutoGen当前进入维护模式，并将新项目引向Agent Framework。本地模型连接使其可以用于内部服务，云端托管产品则提供另一种运行方式。\src{TECH-049}\src{TECH-050}

Agno当前区分SDK、AgentOS运行时与Control Plane管理界面，适合把多个专用助手作为持续服务管理。CrewAI则通过Crews表达协作分工，通过Flows控制整体处理顺序。多Agent适合将资料搜集、数字核对和写作拆成独立工作，再汇总为完整产物；模型调用和结果整合构成相应的运行成本。\src{TECH-051}\src{TECH-052}

Haystack和LlamaIndex更接近资料与检索工作。其中Haystack强调显式组件与流程，LlamaIndex当前公开方向更加重视文档解析及相关产品。开放组件提供可复用的开发能力，商业管理平台和托管服务则承担额外的运行工作。\src{TECH-032}\src{TECH-031}

\subsection{编码Agent与研究人员的文件工作}
OpenCode、pi和Codex CLI提供模型与工具结合的编码工作环境，也可以用于技术研究人员的资料、脚本和文件任务。OpenCode侧重较完整的交互体验，pi侧重精简与可扩展性，Codex CLI提供开源的本地Agent与自定义模型服务配置。\src{TECH-035}\src{TECH-036}\src{TECH-038}

这类工具把脚本、文件和计算环境放在同一工作空间，适合技术研究人员直接处理数据与产物。面向部门同事的工作台则进一步提供浏览器界面、现成文件管理和协作入口。

Claude Code则属于商业产品。当前公开仓库的许可保留所有权利，Anthropic官方也明确不支持通过gateway将其路由到非Claude模型。SGLang等引擎已有相关技术适配文档；软件许可与官方支持范围见附录。\src{TECH-039}\src{TECH-021}

\Needspace{8\baselineskip}
\section{部署形态与完整任务的质量}

\subsection{同一任务的个人与共享参考组合}
个人工作站把模型、资料和工具安排在使用者附近，便于保留个人任务环境；集中服务则在多人之间复用模型资源，并统一更新和管理请求。表\ref{tab:reference-stack}将第\ref{sec:case}节的财报比较任务对应到两套参考设计，说明入口、模型、资料和计算工具如何组成一套系统。

\begin{table}[H]\small\centering
\caption{财报比较任务的个人与共享参考组合}\label{tab:reference-stack}
\begin{tabularx}{\linewidth}{@{}P{22mm}Y Y@{}}\toprule
组成 & 专业个人机：资料与文件助手 & 部门共享：固定财报比较应用\\\midrule
使用入口 & 本机Goose桌面入口，接收目标和文件位置。 & 内网Dify应用，提供文件、期间等固定输入项。\\
模型推理 & 专业Mac运行Gemma 4 31B IT的MLX兼容制品，由MLX-VLM提供本机推理服务。 & GPU服务器上的Qwen3.6-35B-A3B，经vLLM向应用提供推理服务。\\
解析与检索 & 本机配置PDF解析工具，保留页面位置；三份材料可直接读取。 & 应用服务器提取财报和CSV；扩大到资料库时，可另接内网嵌入与检索服务。\\
计算工具 & 本机独立任务目录中的Python计算工具，由Goose调用。 & Dify流程连接隔离计算环境，按规则完成加总、单位及比率检查。\\
结果与状态 & 比较表、说明及来源记录保存在个人任务目录，会话由助手记录。 & 结果文件保存在内部存储，应用数据库记录流程状态；复核与下载入口另行连接。\\
主要建设 & 配置本机模型、解析及计算工具，维护转换制品与任务环境。 & 配置字段、规则、导出与复核，维护身份、数据库、计算环境和共享服务。\\\bottomrule
\end{tabularx}
\source{根据Goose、MLX-VLM、Gemma、Dify及Qwen/vLLM官方能力设计，整体组合为本报告示例。\src{TECH-047}\src{TECH-024}\src{TECH-014}\src{TECH-030}\src{TECH-013}\src{TECH-020}}
\end{table}

个人组合由助手动态安排步骤，适合临时改变比较问题；共享组合提前确定处理规则，便于多位使用者重复形成同类产物。两者也可混合：Goose等个人入口保留本地文件工具，将模型请求交给集中GPU。模型运行位置和工具、资料所在位置可以分别安排。

\subsection{资源、接口与运行条件}
第\ref{sec:memory}节的内存构成可以直接对应到两种部署方式。专业个人机将模型、上下文和本机工具安排在一台设备上；集中GPU服务则通过共享权重和请求调度提高设备利用率。大型模型进入多设备或CPU/GPU异构系统后，主存容量、显存容量与设备互联共同支撑模型运行。

数十名用户可以共用模型服务，各自的数据访问、工具凭证、任务目录和结果按身份管理。这样既可以集中复用推理资源，也可以保留个人文件环境和部门资料权限。模型接口负责传递请求、输出与工具指令，应用负责执行工具、保存来源及恢复任务；具体协议对应关系见附录\ref{app:protocol}。

个人部署的持续成本包含逐台更新和设备利用差异；集中服务可以统一维护并复用资源，同时管理峰值、排队与恢复。容量分析关注实际工作时段和任务组成：注册人数决定服务覆盖范围，同时发起推理的请求数决定瞬时计算负载。

\subsection{从模型速度到业务产物}
财报比较任务把数字、引用、可复算的表格和业务解释放在同一个产物中。程序计算120到138的15\%增长，模型结合财报及注释组织说明；进一步分析增长来源时，再接入销量、价格和产品结构数据。完整任务的质量由这些环节共同形成。

\begin{table}[H]\small\centering
\caption{完整任务的评价层次}
\begin{tabularx}{\linewidth}{@{}P{29mm}Y Y@{}}\toprule
层次 & 核心问题 & 可记录的结果\\\midrule
模型内容 & 事实、数字、事件状态和引用是否正确。 & 准确性、遗漏、冲突、恰当弃答与修订量。\\
推理服务 & 是否在可接受等待内完成请求。 & 首次响应、完整时延、达标吞吐与失败率。\\
应用任务 & 工具和资料能否形成所需产物。 & 任务完成、文件可用性、恢复与人工介入。\\
持续运行 & 用户与运维能否稳定使用。 & 状态恢复、版本变化、资源和维护时间。\\\bottomrule
\end{tabularx}
\end{table}

设备支出决定可用资源，模型调用与维护影响持续成本，人工修订则决定节省时间能否兑现。由此，经济性可以围绕完成一项合格任务所用的机器资源、人员时间和维护投入展开，将推理性能与实际工作效率联系起来。

\Needspace{10\baselineskip}
\section{阶段认识}
当前开放模型已为专业个人工作站和集中服务提供多种能力基础。Qwen与Gemma覆盖长文、多模态和工具使用，稠密与MoE结构提供不同的计算和存储配置；DeepSeek、GLM和Kimi进一步面向复杂研究、办公文件与长过程任务。模型能力、每次计算量和整体内存需求是理解这些路线的三个主要维度。

推理服务把模型能力转化为持续可调用的计算资源。vLLM和SGLang通过调度、缓存与并行服务多人，MLX利用Apple设备的统一内存，异构路线利用CPU和主存扩展模型装载。新模型的专用适配与正式发布包形成不同的更新节奏，对应版本固定、升级和维护的具体工作。

非编码应用则决定同事如何完成任务。财报比较例子中，工作台支持边读边问，Dify将方法固化为重复流程，Goose围绕目标安排文件和计算操作。开发框架进一步连接专用数据与内部流程。这些应用形态已经把模型使用从问答扩展到资料搜集、计算、文件生成和连续任务。

个人环境提供任务连续性和本机操作的便利，共享服务提供资源复用与集中维护，两者可以组合使用。经济性的核心是完成一项工作所消耗的设备资源、人员时间与维护投入：模型提高内容处理能力，服务提高计算利用率，应用减少资料和工具之间的人工衔接，三层共同形成业务效率。

\clearpage
\appendix
\section{版本与资料口径}

\subsection{研究方法与评测口径}
本报告为官方资料研究，尚未执行本项目的本地模型或工作流实验。功能和性能描述注明发布方与任务口径，权重体量取自制品索引；财报比较使用虚构数据，部署组合为架构设计示例。

Gemma的66.4\%与44.1\%对应同卡MRCR v2的128K、8 needle平均结果；GLM的73.0为Toolathlon Verified三次独立运行的pass@1平均结果；Kimi的71.6和54.4来自其模型卡引用的Vals AI金融任务成绩，模型列标为最高思考强度。正文按各自评测解释能力特点，跨模型比较另需统一任务、工具和生成预算。

\subsection{模型制品与上下文}
正文权重体量来自固定模型修订的\texttt{model.safetensors.index.json}中\texttt{metadata.total\_size}，按十进制GB换算并保留两位小数。该数字统计对应制品的数值数组（tensor）字节；运行内存另包含上下文缓存与计算开销。

Qwen卡片的语言模型参数口径与完整多模态制品不同。Gemma的API元素统计与分片索引中另报的参数数存在差异，本文保留原始记录，并使用索引字节数描述体量。DeepSeek与Kimi使用混合量化及打包数组，逻辑参数量采用模型卡口径，制品大小采用实际字节口径。

DeepSeek表中的条件记忆指Engram，通过按输入查找的方式稀疏访问相应参数。Kimi采用MXFP4等混合低精度格式；FP8、BF16分别涉及8位和16位浮点表示，混合制品按各部分的实际字节合计。

上下文数字来自模型卡或配置，表示对应资料中的能力或配置上限。运行时可用长度还受内存、并发、引擎与具体任务质量影响。全部8项模型的原始索引及模型修订均列入报告来源清单。

\subsection{后续模型观察}
Meta已公开讨论Muse Spark 1.2的开放权重计划，本轮资料取得的是计划信息。Qwen4、GLM-5.4、DeepSeek-V4.1-Pro及Kimi-K3.1等后续型号列入持续观察，公开权重、许可与规格发布后可沿本报告的模型、推理和应用三层结构补充分析。\src{TECH-018}

\subsection{软件版本快照}
下表列示本轮取得的部分版本信息。源码HEAD、GitHub latest-release、包索引和在线文档分别记录修订与更新时间，部署支持表据此区分正式版本和后续适配。

{\small
\begin{longtable}{@{}P{44mm}P{41mm}P{69mm}@{}}
\caption{主要软件的版本与成熟度资料}\label{tab:versions}\\
\toprule 项目 & 本轮版本记录 & 资料说明\\\midrule
\endfirsthead
\toprule 项目 & 本轮版本记录 & 资料说明\\\midrule
\endhead
\bottomrule\endfoot
vLLM & v0.29.0 & 9月18日再次读取latest-release，仍为该版本。\\
SGLang & v0.5.19 & 9月18日再次读取latest-release，仍为该版本。\\
TokenSpeed & v0.1.0 & 7月24日发布；本文模型配置来自后续源码。\\
MLX-LM & 0.31.3 & GitHub/PyPI快照；所读源码另有固定修订。\\
MLX-VLM & v0.7.1 & 源码与文档能力仍按对应修订判断。\\
llama.cpp & v0.4.1 & 模型格式、后端和API由具体构建决定。\\
KTransformers & v0.7.1 & kt-kernel及服务组合需要对应依赖。\\
ik\_llama.cpp & 固定源码commit & 本轮latest-release API未取得记录。\\
Open WebUI & v0.11.3 & 产品功能与品牌许可采用当前资料快照。\\
LibreChat & 源码及当前文档 & latest-release API未取得记录，部分新功能为实验性。\\
Onyx & v4.7.7 & Community与Enterprise内容区分。\\
RAGFlow & v0.27.2 & 资料处理及代码沙箱分别配置。\\
Hermes Agent & v2026.9.14 & 桌面与Skills文档同时保存源码及正式标签。\\
LobeHub & v2.2.17 & 开发分支canary；Codex整合另存正式标签文档。\\
Cherry Studio & v2.0.14 & Agent运行时与Code Mate文档另存正式标签。\\
OpenClaw & v2026.9.4 & 团队与Codex整合另存正式标签文档。\\
Agno & v3.0.10 & SDK、运行时与管理产品分层。\\
CrewAI & 1.15.22 & 已取得对应版本概念文档。\\
AnythingLLM & v1.16.1 & 桌面、Docker及商业产品区分。\\
Microsoft Agent Framework & .NET子包1.21.0 & 不是所有语言包的统一版本号。\\
Codex CLI & 正式release 0.155.0 & 9月18日取得新发布记录及官方执行文档。\\
Claude Code & v2.1.276 & CLI与Agent SDK能力分别记录。\\
\end{longtable}
}
\smallnote{依据各官方repo的release记录、所读包索引及专题文档；完整修订、时间与哈希保存在报告数据清单。}

\subsection{自定义模型许可的门槛与定义}\label{app:model-license}
GLM-5.3许可第2项同时考察两个条件：被许可方或关联方经营模型即服务（Model as a Service，MaaS）业务，以及被许可方及关联方在任意连续12个月内合计总收入超过100亿美元。两项均成立时，任何商业使用前需通过Z.AI安全审查。\src{TECH-015}（LICENSE第2项）

该许可中的MaaS强调第三方对推理或微调输入、参数或训练数据的实质控制，并排除仅将模型嵌入特定功能的终端产品及单纯转发其他方托管请求的情况。触发条件同时包括相应服务业务与集团合计收入。

Kimi K3的第2项将相应集团收入门槛设为连续12个月2,000万美元，并在同时经营MaaS时要求单独协议；第3项对使用模型的商业产品或服务规定署名条件，其门槛为超过1亿月活或每月2,000万美元收入。第4项则规定，符合定义的内部使用，以及通过Moonshot官方产品或认证推理合作方的使用，不适用第2、3项。\src{TECH-017}（LICENSE第2--4项）

Kimi的内部使用定义要求软件、输出及底层能力均不向第三方提供。声明保留等其余许可义务继续适用。

\subsection{软件许可与产品边界}\label{app:software-license}
\begin{table}[H]\small\centering
\caption{软件核心与商业产品的许可区别}
\begin{tabularx}{\linewidth}{@{}P{35mm}Y@{}}\toprule
类别 & 本轮资料所体现的条件\\\midrule
标准许可核心 & 多数引擎及开发框架采用MIT或Apache 2.0，第三方组件、模型和企业支持另行适用。\\
Open WebUI & LICENSE第4项规定品牌保留；滚动30天内不超过50名自然人、书面许可或企业授权等为相应例外。该人数用于判断品牌修改例外。\\
LobeHub & Community License基于Apache并附加商用条件；原版可商用，开发并分发衍生作品需商业许可。\\
Cherry Studio & 社区版AGPL-3.0，允许商用；修改、分发及条款规定的网络交互涉及相应源码义务。企业版另提供集中管理与商业授权。\\
Hermes / OpenClaw & 核心均为MIT，模型、第三方扩展及工具服务各自适用相应授权。\\
Dify & 修改版Apache许可对多租户及前端标识设限；每个workspace计为一个租户。多租户服务和去除或修改前端标识涉及附加授权；不使用其前端时，标识限制不适用。\\
Onyx / Flowise & 普通代码与指定企业目录或文件分别授权。\\
n8n & Sustainable Use及企业许可，内部用途与对外提供方式分别判断。\\
Agno等产品体系 & 开放SDK与商业管理界面、托管产品分别授权。\\
Claude Code & 当前repo保留所有权利，软件修改与使用适用其商业授权。\\\bottomrule
\end{tabularx}
\source{对应项目LICENSE及产品说明。\src{TECH-043}\src{TECH-057}\src{TECH-058}\src{TECH-056}\src{TECH-061}\src{TECH-030}\src{TECH-045}\src{TECH-033}\src{TECH-034}\src{TECH-051}\src{TECH-039}}
\end{table}

\subsection{应用连接的协议差异}\label{app:protocol}
Chat Completions、Responses和Anthropic Messages采用不同请求与事件形式，工具调用、思考内容、状态和流式输出的处理也有差别。下表将已归档接口与相应客户端配置对应。

\begin{table}[H]\small\centering
\caption{本轮核对的部分应用接口差异}
\begin{tabularx}{\linewidth}{@{}P{34mm}Y Y@{}}\toprule
组件 & 已核对的接口信息 & 对应用连接的影响\\\midrule
MLX-LM服务快照 & 原生POST路由为Completions及Chat Completions。 & 对应采用Completions及Chat Completions的客户端。\\
MLX-VLM / llama.cpp & 当前文档列出Responses；llama.cpp另列Messages。 & 提供多协议接入，工具循环配套相应模型解析器。\\
Codex CLI & 当前自定义模型服务配置使用Responses。 & 服务需要满足该协议及实际请求所用功能。\\
OpenCode / pi & 提供自定义模型连接；pi可明确选择协议。 & 模型、工具解析与上下文设置共同影响任务表现。\\\bottomrule
\end{tabularx}
\source{固定源码或官方配置文档。\src{TECH-023}\src{TECH-024}\src{TECH-025}\src{TECH-038}\src{TECH-035}\src{TECH-036}}
\end{table}

\clearpage
\section{参考文献}
本文引用对应2026年9月16日至18日取得的官方资料快照。模型及软件的具体版本、原件与内容校验记录另附于报告源码目录。正文的产品功能来自文档说明，业务例子和部署分析为本报告的归纳。
\referenceentry{TECH-020}{vLLM：项目与模型服务文档}{vllm-project}{资料快照 2026-09-17，repo修订 6ca2b23e22aa}{https://github.com/vllm-project/vllm}
\referenceentry{TECH-029}{LangGraph：状态流程、持久运行与人工中断}{langchain-ai}{资料快照 2026-09-17，repo修订 230927fb3a9a}{https://github.com/langchain-ai/langgraph}
\referenceentry{TECH-037}{DeepSeek Harness：架构、模型与预览状态}{deepseek-ai}{资料快照 2026-09-17，repo修订 0d1f50007f9b}{https://github.com/deepseek-ai/deepseek-harness}
\referenceentry{TECH-004}{Qwen3.8-27B：模型卡、配置与许可}{Qwen}{资料快照 2026-09-16，repo修订 1d4bf0f2ff60}{https://huggingface.co/Qwen/Qwen3.8-27B}
\referenceentry{TECH-013}{Qwen3.6-35B-A3B：模型卡、配置与许可}{Qwen}{资料快照 2026-09-16，repo修订 995ad96eacd9}{https://huggingface.co/Qwen/Qwen3.6-35B-A3B}
\referenceentry{TECH-014}{Gemma 4 31B：模型卡与配置}{google}{资料快照 2026-09-16，repo修订 842da3794eaa}{https://huggingface.co/google/gemma-4-31B-it}
\referenceentry{TECH-006}{Gemma 4 26B A4B：模型卡与配置}{google}{资料快照 2026-09-16，repo修订 4d7ae4984b7d}{https://huggingface.co/google/gemma-4-26B-A4B-it}
\referenceentry{TECH-008}{DeepSeek-V4.1-Flash：模型卡、配置与许可}{deepseek-ai}{资料快照 2026-09-16，repo修订 dba1be0a40aa}{https://huggingface.co/deepseek-ai/DeepSeek-V4.1-Flash}
\referenceentry{TECH-016}{GLM-5.3-Flash：模型卡与MIT许可}{zai-org}{资料快照 2026-09-16，repo修订 eb9eb208eb0d}{https://huggingface.co/zai-org/GLM-5.3-Flash}
\referenceentry{TECH-015}{GLM-5.3：模型卡与自定义许可}{zai-org}{资料快照 2026-09-16，repo修订 aca966e4e027}{https://huggingface.co/zai-org/GLM-5.3}
\referenceentry{TECH-017}{Kimi K3：模型卡与自定义许可}{moonshotai}{资料快照 2026-09-16，repo修订 f831ab668142}{https://huggingface.co/moonshotai/Kimi-K3}
\referenceentry{TECH-012}{Gemma 4：Apache 2.0许可}{Google}{资料快照 2026-09-16}{https://ai.google.dev/gemma/docs/gemma_4_license}
\referenceentry{TECH-019}{Apache License 2.0}{Apache Software Foundation}{资料快照 2026-09-16}{https://www.apache.org/licenses/LICENSE-2.0.txt}
\referenceentry{TECH-021}{SGLang：项目与服务接口文档}{sgl-project}{资料快照 2026-09-17，repo修订 b02e16a895ad}{https://github.com/sgl-project/sglang}
\referenceentry{TECH-022}{TokenSpeed：调度、模型与服务文档}{lightseekorg}{资料快照 2026-09-17，repo修订 688bbb04f5b4}{https://github.com/lightseekorg/tokenspeed}
\referenceentry{TECH-040}{TensorRT-LLM：项目与许可}{NVIDIA}{资料快照 2026-09-17，repo修订 fffdfba7279a}{https://github.com/NVIDIA/TensorRT-LLM}
\referenceentry{TECH-026}{KTransformers：异构推理与GLM Flash教程}{kvcache-ai}{资料快照 2026-09-17，repo修订 0b954fe4ca2f}{https://github.com/kvcache-ai/ktransformers}
\referenceentry{TECH-009}{vLLM v0.29.0：固定正式版架构注册表}{vLLM project}{资料快照 2026-09-16}{https://github.com/vllm-project/vllm/releases/tag/v0.29.0}
\referenceentry{TECH-010}{Qwen3.8-27B：vLLM官方部署配置与上游测试构建}{vLLM project}{资料快照 2026-09-16}{https://recipes.vllm.ai/Qwen/Qwen3.8-27B}
\referenceentry{TECH-011}{DeepSeek-V4.1-Flash：vLLM初始支持合并记录}{vLLM project}{资料快照 2026-09-16}{https://github.com/vllm-project/vllm/issues/56400}
\referenceentry{TECH-054}{vLLM v0.29.0：发布记录与服务benchmark参数}{vllm-project}{资料快照 2026-09-18}{https://docs.vllm.ai/en/v0.29.0/cli/bench/serve/}
\referenceentry{TECH-023}{MLX-LM：项目与服务实现}{ml-explore}{资料快照 2026-09-17，repo修订 872ae88d1fac}{https://github.com/ml-explore/mlx-lm}
\referenceentry{TECH-042}{Apple WWDC26：在Mac上运行本地Agent}{Apple}{资料快照 2026-09-17}{https://developer.apple.com/videos/play/wwdc2026/232/}
\referenceentry{TECH-024}{MLX-VLM：多模态与服务文档}{Blaizzy}{资料快照 2026-09-17，repo修订 c72f637e0378}{https://github.com/Blaizzy/mlx-vlm}
\referenceentry{TECH-025}{llama.cpp：项目与HTTP服务文档}{ggml-org}{资料快照 2026-09-17，repo修订 c6824a9e42ce}{https://github.com/ggml-org/llama.cpp}
\referenceentry{TECH-027}{ik\_llama.cpp：量化与工具调用文档}{ikawrakow}{资料快照 2026-09-17，repo修订 e64c0ea59f75}{https://github.com/ikawrakow/ik_llama.cpp}
\referenceentry{TECH-055}{SGLang：吞吐与内存调优指南}{sgl-project}{资料快照 2026-09-18}{https://docs.sglang.io/docs/advanced_features/hyperparameter_tuning.md}
\referenceentry{TECH-043}{Open WebUI：产品功能与许可}{open-webui}{资料快照 2026-09-18，repo修订 0a7c15832fb3}{https://github.com/open-webui/open-webui}
\referenceentry{TECH-044}{LibreChat：Agents与代码执行服务}{danny-avila}{资料快照 2026-09-18，repo修订 12d78909d3f5}{https://github.com/danny-avila/LibreChat}
\referenceentry{TECH-045}{Onyx：研究平台与产品版本}{onyx-dot-app}{资料快照 2026-09-18，repo修订 d7cc16686a35}{https://github.com/onyx-dot-app/onyx}
\referenceentry{TECH-030}{Dify：产品、部署与许可}{langgenius}{资料快照 2026-09-17，repo修订 9c6c48b50ba9}{https://github.com/langgenius/dify}
\referenceentry{TECH-056}{Hermes Agent：桌面、记忆与Skills}{NousResearch}{资料快照 2026-09-18，repo修订 e83b1d51f13a}{https://github.com/NousResearch/hermes-agent}
\referenceentry{TECH-049}{Microsoft Agent Framework：项目与BUILD 2026更新}{microsoft}{资料快照 2026-09-18，repo修订 459fe6eca832}{https://github.com/microsoft/agent-framework}
\referenceentry{TECH-057}{LobeHub：Agent工作空间、自托管与Coding Agent整合}{lobehub}{资料快照 2026-09-18，repo修订 270a4a50cb32}{https://github.com/lobehub/lobehub}
\referenceentry{TECH-058}{Cherry Studio：桌面、Agent运行时与企业版}{CherryHQ}{资料快照 2026-09-18，repo修订 142f3463f287}{https://github.com/CherryHQ/cherry-studio}
\referenceentry{TECH-053}{AnythingLLM：个人与团队资料应用}{Mintplex-Labs}{资料快照 2026-09-18，repo修订 90108f98f295}{https://github.com/Mintplex-Labs/anything-llm}
\referenceentry{TECH-046}{RAGFlow：文档与Agent能力}{infiniflow}{资料快照 2026-09-18，repo修订 15b6668dd8f7}{https://github.com/infiniflow/ragflow}
\referenceentry{TECH-034}{n8n：自动化平台与许可}{n8n-io}{资料快照 2026-09-17，repo修订 f960330ecb34}{https://github.com/n8n-io/n8n}
\referenceentry{TECH-033}{Flowise：产品及分层许可}{FlowiseAI}{资料快照 2026-09-17，repo修订 9291856d1ea4}{https://github.com/FlowiseAI/Flowise}
\referenceentry{TECH-061}{OpenClaw：团队、记忆与Codex Harness整合}{openclaw}{资料快照 2026-09-18，repo修订 cc96aefe25bb}{https://github.com/openclaw/openclaw}
\referenceentry{TECH-060}{Claude Code与Agent SDK：通用任务、Skills与MCP}{anthropics}{资料快照 2026-09-18，repo修订 31a3b00bef14}{https://github.com/anthropics/claude-code}
\referenceentry{TECH-059}{Codex：Skills、MCP、非交互执行与SDK}{openai}{资料快照 2026-09-18，repo修订 7498521d288b}{https://github.com/openai/codex}
\referenceentry{TECH-062}{Agent Skills：业务方法、脚本与资源格式}{Agent Skills}{资料快照 2026-09-18}{https://agentskills.io/home}
\referenceentry{TECH-063}{Model Context Protocol：资料、工具与系统连接}{Model Context Protocol}{资料快照 2026-09-18}{https://modelcontextprotocol.io/docs/2026-07-28/getting-started/intro}
\referenceentry{TECH-035}{OpenCode：模型连接与权限}{anomalyco}{资料快照 2026-09-17，repo修订 2e018f70f2d1}{https://github.com/anomalyco/opencode}
\referenceentry{TECH-036}{pi：Agent、模型与执行环境}{earendil-works}{资料快照 2026-09-17，repo修订 509ee2bd0ba9}{https://github.com/earendil-works/pi}
\referenceentry{TECH-038}{Codex CLI：项目许可与官方配置文档}{openai}{资料快照 2026-09-17，repo修订 fd346b8dbaa2}{https://github.com/openai/codex}
\referenceentry{TECH-039}{Claude Code：许可、网关与企业部署文档}{anthropics}{资料快照 2026-09-17，repo修订 b782847db9a1}{https://github.com/anthropics/claude-code}
\referenceentry{TECH-051}{Agno：SDK、AgentOS与管理产品}{agno-agi}{资料快照 2026-09-18，repo修订 a71c4d0c6072}{https://github.com/agno-agi/agno}
\referenceentry{TECH-052}{CrewAI：Crews、Flows与模型连接}{crewAIInc}{资料快照 2026-09-18，repo修订 597b99fb576b}{https://github.com/crewAIInc/crewAI}
\referenceentry{TECH-032}{Haystack：检索与Agent流程}{deepset-ai}{资料快照 2026-09-17，repo修订 0defdcff6495}{https://github.com/deepset-ai/haystack}
\referenceentry{TECH-031}{LlamaIndex：项目与文档处理方向}{run-llama}{资料快照 2026-09-17，repo修订 fd4a517ad649}{https://github.com/run-llama/llama_index}
\referenceentry{TECH-028}{LangChain：框架概览}{langchain-ai}{资料快照 2026-09-17，repo修订 5c1f28271295}{https://github.com/langchain-ai/langchain}
\referenceentry{TECH-041}{Deep Agents：Agent能力与本地模型连接}{langchain-ai}{资料快照 2026-09-17，repo修订 26d2eae2ec07}{https://github.com/langchain-ai/deepagents}
\referenceentry{TECH-050}{AutoGen：维护状态说明}{microsoft}{资料快照 2026-09-18，repo修订 027ecf0a379b}{https://github.com/microsoft/autogen}
\referenceentry{TECH-018}{Muse Spark 1.2：多模态能力与开放权重计划}{Meta}{资料快照 2026-09-16}{https://research.meta.ai/blog/multimodal-intelligence-of-muse-spark-1-2}

\end{document}
